\chapter{Vector Spaces and Subspaces}

The following table holds three records of network traffic; each row is one record, each column
one feature (incoming packets, outgoing packets, duration, failed connections). Read from top to
bottom, the columns are
\[
  \mathbf{c}_{1}=(1,2,3),\quad \mathbf{c}_{2}=(2,4,6),\quad
  \mathbf{c}_{3}=(0,1,2),\quad \mathbf{c}_{4}=(1,3,5),
\]
while the rows are $\mathbf{r}_{1}=(1,2,0,1)$, $\mathbf{r}_{2}=(2,4,1,3)$, and
$\mathbf{r}_{3}=(3,6,2,5)$.

How many columns actually carry new information? Counting the nonzero columns gives four, and
that count is wrong: the second column is twice the first, the fourth is the sum of the first
and the third. The table holds only two independent directions of data. The sentence ``can be
recomputed'' acquires meaning only once we can measure how few vectors build it up.

Chapters 1 through 5 prepared the language of sets, the axioms of the real numbers, and then
limits of sequences and series; all of them work with closeness. The problem above contains no
closeness at all: its entries are exact, and the question is purely algebraic. Which sets are
closed under both operations, and how few vectors build them. One habit from Chapter 5 stays in
use: here every combination involves finitely many terms.

There is a second question that demands the same tool. A linear scorer assigns the score
$s(\mathbf{x})=\mathbf{w}\cdot\mathbf{x}$ with $\mathbf{w}=(1,1,1)$. A direction
$\boldsymbol{\delta}$ that does not change the score satisfies
$\mathbf{w}\cdot\boldsymbol{\delta}=0$, which is a plane; every perturbation splits into a part
that changes the score and a part that does not. Is that split unique? Yes, and the property is
called a direct sum.

\begin{untukanak}
Picture a box of poster paint holding four cans: red, yellow, blue, and orange. We may mix the
contents of the cans in any proportion, and the mixture may itself be used to mix again. Orange
can be made from red and yellow. So the orange can adds not a single capability that the box did
not already have; it only takes up space.

The nagging question: how few cans must remain so that every color we can make stays a color we
can make? The answer is not the same as the number of cans left after we throw one away, and the
order of removal must not change that number. That is no accident, and proving it is not as
simple as it looks.

There is a stranger question. Suppose two people mix according to different recipes, and the
results are exactly the same. Are the recipes the same? Sometimes yes, sometimes no. In this
chapter we will learn why that number has a name of its own, why it does not depend on the order
of removal, and why two different mixtures sometimes yield the same color.
\end{untukanak}

\section{The vector space axioms and their examples}
\tingkat{basic}

A vector space is a set whose members may be added and may be enlarged or shrunk by a number,
with the result always staying inside that set, and with the same rules of computation as for
ordinary pairs of numbers.

\begin{definisi}[Vector space]
\index{vector space}
A nonempty set $V$ together with an addition $V\times V\to V$ and a scalar multiplication
$\mathbb{R}\times V\to V$ is called a \emph{vector space over $\mathbb{R}$} if for all
$\mathbf{u},\mathbf{v},\mathbf{w}\in V$ and all $a,b\in\mathbb{R}$ the following hold:
\begin{enumerate}
  \item $\mathbf{u}+\mathbf{v}=\mathbf{v}+\mathbf{u}$;
  \item $(\mathbf{u}+\mathbf{v})+\mathbf{w}=\mathbf{u}+(\mathbf{v}+\mathbf{w})$;
  \item there is $\mathbf{0}\in V$ with $\mathbf{v}+\mathbf{0}=\mathbf{v}$ for every
        $\mathbf{v}$;
  \item there is $-\mathbf{v}\in V$ with $\mathbf{v}+(-\mathbf{v})=\mathbf{0}$;
  \item $a(\mathbf{u}+\mathbf{v})=a\mathbf{u}+a\mathbf{v}$;
  \item $(a+b)\mathbf{v}=a\mathbf{v}+b\mathbf{v}$;
  \item $a(b\mathbf{v})=(ab)\mathbf{v}$;
  \item $1\mathbf{v}=\mathbf{v}$.
\end{enumerate}
\end{definisi}

Both operations carry the closure requirement inside their type: the result always lies in $V$,
so this definition contains no separate closure axiom. The eight axioms are not independent of
one another either, and we do not chase the shortest list; what we need are their small
consequences, all derived without looking at the shape of the members of $V$.

\begin{lemma}
In any vector space, the zero vector is unique, the additive inverse of every vector is unique,
$0\mathbf{v}=\mathbf{0}$, $(-1)\mathbf{v}=-\mathbf{v}$, and $c\mathbf{v}=\mathbf{0}$ implies
$c=0$ or $\mathbf{v}=\mathbf{0}$.
\end{lemma}

\par\medskip\noindent\textbf{Proof.}\ 
For uniqueness of zero, suppose $\mathbf{0}$ and $\mathbf{0}'$ are both neutral. Then
$\mathbf{0}=\mathbf{0}+\mathbf{0}'=\mathbf{0}'$. For uniqueness of inverses, suppose
$\mathbf{v}+\mathbf{a}=\mathbf{0}$ and $\mathbf{v}+\mathbf{b}=\mathbf{0}$. Then
\[
  \mathbf{a}=\mathbf{a}+\mathbf{0}=\mathbf{a}+(\mathbf{v}+\mathbf{b})
  =(\mathbf{a}+\mathbf{v})+\mathbf{b}=\mathbf{0}+\mathbf{b}=\mathbf{b}.
\]
Next $0\mathbf{v}=(0+0)\mathbf{v}=0\mathbf{v}+0\mathbf{v}$. Adding $-(0\mathbf{v})$ to both
sides gives $\mathbf{0}=0\mathbf{v}$. From there
$\mathbf{v}+(-1)\mathbf{v}=1\mathbf{v}+(-1)\mathbf{v}=(1-1)\mathbf{v}=0\mathbf{v}=\mathbf{0}$,
so $(-1)\mathbf{v}$ is an inverse of $\mathbf{v}$, and since inverses are unique,
$(-1)\mathbf{v}=-\mathbf{v}$. Finally, if $c\mathbf{v}=\mathbf{0}$ and $c\neq 0$, then
\[
  \mathbf{v}=1\mathbf{v}=(c^{-1}c)\mathbf{v}=c^{-1}(c\mathbf{v})=c^{-1}\mathbf{0}=\mathbf{0}.
\]
\hfill$\square$\par\medskip

The following four examples suffice for this entire book.

\emph{(i) Coordinate space.} $\mathbb{R}^{n}$ with addition and scalar multiplication coordinate
by coordinate. Here $\mathbf{0}=(0,\dots,0)$, and the eight axioms reduce to the field properties
of the real numbers already discussed in Chapter 3.

\emph{(ii) Matrix space.} $M_{2\times3}(\mathbb{R})$, the set of $2\times3$ matrices with real
entries, with operations entry by entry and the zero matrix as $\mathbf{0}$. Nothing changes if
$2\times3$ is replaced by $m\times n$.

\emph{(iii) Polynomial space.} $P_{3}$ is the set of polynomials with real coefficients of degree
at most $3$, together with the zero polynomial, with operations term by term. A trap worth
noting: polynomials of degree \emph{exactly} $3$ do not form a vector space, since $x^{3}$ and
$-x^{3}+x$ both have degree three but their sum $x$ has degree one. The words ``at most'' are not
decoration.

\emph{(iv) Function space.} The set of all functions $\mathbb{R}\to\mathbb{R}$, written
$\mathbb{R}^{\mathbb{R}}$, with $(f+g)(x)=f(x)+g(x)$ and $(af)(x)=af(x)$. The zero function is
$\mathbf{0}$, and the vector space axioms are inherited from $\mathbb{R}$ at each point $x$.

Example (iv) yields a subtler example. The set $C[0,1]$ of continuous functions on $[0,1]$ is a
vector space, because sums and scalar multiples of continuous functions remain continuous. So is
the set of convergent real sequences, using the limit-of-sums property from Chapter 4. Both are
parts of a larger space, and parts of that kind we now define.

Back to $\mathbb{R}^{2}$, where both operations are visible directly.
Figure~\ref{fig:ch06-dua-operasi} uses $\mathbf{u}=(2,1)$ and $\mathbf{v}=(1,2)$: the sum
$\mathbf{u}+\mathbf{v}=(3,3)$ is reached along two orders, and $2\mathbf{u}$ as well as
$-\mathbf{u}$ stay collinear with $\mathbf{u}$.

\begin{figure}[htbp]
\centering
\begin{tikzpicture}[scale=0.85,>=Stealth,font=\small]
  \draw[->,layergrey] (-2.8,0) -- (5.0,0) node[below] {$x$};
  \draw[->,layergrey] (0,-1.8) -- (0,4.2) node[left] {$y$};
  \draw[dashed,layergrey] (2,1) -- (3,3);
  \draw[dashed,layergrey] (1,2) -- (3,3);
  \draw[->,thick,accent] (0,0) -- (2,1) node[below right] {$\mathbf{u}$};
  \draw[->,thick,layerindigo] (0,0) -- (1,2) node[above left] {$\mathbf{v}$};
  \draw[->,thick,layerteal] (0,0) -- (3,3) node[right] {$\mathbf{u}+\mathbf{v}$};
  \draw[->,layerorange] (0,0) -- (4,2) node[right] {$2\mathbf{u}$};
  \draw[->,layerbrown] (0,0) -- (-2,-1) node[left] {$-\mathbf{u}$};
  \fill[ink] (0,0) circle (1.2pt);
  \node[below left] at (0,0) {$\mathbf{0}$};
\end{tikzpicture}
\caption{The two operations in $\mathbb{R}^{2}$. The thick arrows show $\mathbf{u}=(2,1)$ and
$\mathbf{v}=(1,2)$; the dashed lines complete the parallelogram, so that
$\mathbf{u}+\mathbf{v}=(3,3)$ is reached along two identical orders. The thin arrows show
$2\mathbf{u}=(4,2)$ and $-\mathbf{u}=(-2,-1)$, both staying on one line with
$\mathbf{u}$.}
\label{fig:ch06-dua-operasi}
\end{figure}

\kamus{vector space}{ruang vektor}{a mixing box: its contents may be added and scaled, and the result never leaves the box.}
\begin{bacakeras}
Eight axioms, read as a checklist. Two voices, alternating; the second voice always answers in
three words.

First voice: Swap the places. Second voice: The result is the same.
First voice: Move the brackets. Second voice: The result is the same.
First voice: There is a zero vector. Second voice: It adds nothing.
First voice: There is an opposite vector. Second voice: Back to zero.
First voice: A scalar spreads over two terms. Second voice: Left and right agree.
First voice: Two scalars add up. Second voice: One vector carries it.
First voice: Two scalars multiply. Second voice: One multiplication only.
First voice: The scalar one. Second voice: Silent, unchanged.

The last two lines sound almost alike although they differ: one adds proportions, the other
multiplies proportions. That is where the most frequent arithmetic error is born.

Say it once more, slowly, in a single breath: swap the places, move the brackets, there is a
zero, there is an opposite, it spreads, it adds up, it multiplies, it is silent.
\end{bacakeras}


\section{Subspaces and the subspace test}
\tingkat{basic}

A subspace is a part of a vector space that still behaves as a vector space even when we use only
the members of that part.

\begin{definisi}[Subspace]
\index{subspace}
A set $W\subseteq V$ is called a \emph{subspace} of $V$ if $W$ is a vector space with the same
operations as on $V$.
\end{definisi}

That definition demands eight axioms, and checking eight things every time is too expensive. Six
of them are free: commutativity, associativity, distributivity hold in $W$ because they hold in
$V$. Three things remain, and they are the content of the next theorem.

If $W$ is nonempty and closed under both operations, then $W$ is automatically a subspace.

\begin{teorema}[Subspace test]
\index{subspace test}
Let $V$ be a vector space and $W\subseteq V$ with $W\neq\varnothing$. Then $W$ is a subspace of
$V$ if and only if $W$ is closed under addition and scalar multiplication.
\end{teorema}

\par\medskip\noindent\textbf{Proof.}\ 
If $W$ is a subspace, the result of both operations on members of $W$ must be a member of $W$, so
both closures hold. Conversely, suppose $W$ is nonempty and closed under both operations. Take
$\mathbf{w}\in W$. By scalar closure, $\mathbf{0}=0\mathbf{w}\in W$, so $W$ contains the neutral
element. Also $-\mathbf{w}=(-1)\mathbf{w}\in W$, so every member of $W$ has its inverse inside
$W$. The remaining axioms hold because they hold in $V$ and the operations on $W$ are merely the
restrictions of the operations on $V$.
\hfill$\square$\par\medskip

In practice we check three things: $\mathbf{0}\in W$ and the two closures. All three are
necessary; the next three examples show that failure of one condition cannot be covered up by
another.
\begin{resep}
\begin{enumerate}
  \item Write the membership condition of $W$ as an equation, not as a picture or a name.
  \item Substitute $\mathbf{v}=\mathbf{0}$ into that condition. If it fails, stop: not a subspace.
  \item Take arbitrary $\mathbf{u},\mathbf{v}\in W$, check $\mathbf{u}+\mathbf{v}$ using the same
        condition, not using numerical examples.
  \item Take arbitrary $c\in\mathbb{R}$ and $\mathbf{u}\in W$, check $c\mathbf{u}$; include
        $c<0$ and $0<c<1$.
\end{enumerate}
A warning about order: do not swap step 2 with steps 3 and 4. A set that does not contain
$\mathbf{0}$ need not be tested for closure, but stopping after step 2 is also wrong: containing
$\mathbf{0}$ is not enough, as the following examples show.
\end{resep}


\emph{A line not through the origin.} $W=\{(x,y)\in\mathbb{R}^{2}:x+y=1\}$ does not contain
$\mathbf{0}$, so it is not a subspace; it is also closed under neither addition nor scalar
multiplication, since $(1,0)+(0,1)$ and $2(1,0)$ both leave $W$.

\emph{The union of two axes.} $W=\{(x,0)\}\cup\{(0,y)\}$ contains $\mathbf{0}$ and is closed
under scalar multiplication, but $(1,0)+(0,1)=(1,1)$ leaves $W$.

\emph{The first quadrant.} $W=\{(x,y):x\ge 0,\ y\ge 0\}$ is closed under addition, but
$(-1)(1,1)=(-1,-1)$ leaves $W$.

All three examples fail a different condition; Figure~\ref{fig:ch06-garis} shows the simplest
pair, a line through the origin and a line that misses it.

\begin{figure}[htbp]
\centering
\begin{tikzpicture}[scale=0.72,>=Stealth,font=\small]
  \draw[->,layergrey] (-1.6,0) -- (3.2,0) node[below] {$x$};
  \draw[->,layergrey] (0,-1.6) -- (0,3.2) node[left] {$y$};
  \draw[very thick,layerteal] (-1.4,-1.4) -- (2.8,2.8);
  \fill[ink] (0,0) circle (1.3pt) node[below left] {$\mathbf{0}$};
  \fill[accent] (2,2) circle (1.5pt) node[above right] {$(2,2)$};
  \node at (0.8,-2.4) {\emph{a line through the origin}};
  \begin{scope}[xshift=7cm]
    \draw[->,layergrey] (-1.6,0) -- (3.2,0) node[below] {$x$};
    \draw[->,layergrey] (0,-1.6) -- (0,3.2) node[left] {$y$};
    \draw[very thick,accent] (-0.8,1.8) -- (2.4,-1.4);
    \fill[accent] (1,0) circle (1.5pt) node[below right] {$(1,0)$};
    \fill[accent] (0,1) circle (1.5pt) node[above left] {$(0,1)$};
    \draw[ink,thick] (0,0) circle (3pt);
    \node[below left] at (0,0) {$\mathbf{0}$};
    \node at (0.8,-2.4) {\emph{a line without the origin}};
  \end{scope}
\end{tikzpicture}
\caption{Two sets in $\mathbb{R}^{2}$. Left: $\{(x,y):y=x\}$ contains $\mathbf{0}$ and
$(2,2)$, so it is closed under both operations. Right: $\{(x,y):x+y=1\}$ contains $(1,0)$ and
$(0,1)$, but the origin is marked with an empty circle because it is not included; a single
missing point is enough to disqualify a subspace.}
\label{fig:ch06-garis}
\end{figure}



Conversely, the next two classes of examples always pass.

\emph{Null space.} For an $m\times n$ matrix $A$, the set
$\{\mathbf{x}:A\mathbf{x}=\mathbf{0}\}$ is a subspace, since $A\mathbf{0}=\mathbf{0}$,
$A(\mathbf{x}+\mathbf{y})=A\mathbf{x}+A\mathbf{y}=\mathbf{0}$, and
$A(c\mathbf{x})=cA\mathbf{x}=\mathbf{0}$. That set is called the null space or kernel of $A$,
written $\ker A$, and is discussed further in Chapter 7.

\emph{Polynomials vanishing at one point.} The set
$\{p\in P_{3}:p(1)=0\}$ is a subspace of $P_{3}$, because evaluation at the point $1$ is linear:
$(p+q)(1)=p(1)+q(1)$ and $(cp)(1)=c\,p(1)$.

Both examples preserve the solution set of a homogeneous equation. That is no accident.
\begin{tantangan}
The first-quadrant example fails because of one negative scalar. Find a set
$W\subseteq\mathbb{R}^{2}$ that contains $\mathbf{0}$, is closed under addition, is not closed
under scalar multiplication, and is neither the union of two axes, nor a quadrant, nor a set with
integer coordinates.

Hint 1: do not regulate the direction of the arrows; regulate which numbers may appear in the sum
of the two coordinates.

Hint 2: require $x+y$ to be an integer. Then $(1,0)+(0,1)=(1,1)$ stays inside, whereas
$\tfrac{1}{2}(1,0)$ leaves.

If you get stuck, go back to the section ``Subspaces and the subspace test'' and compare with the
line example $x+y=1$; the only difference is the right-hand side of the equation.
\end{tantangan}


\kamus{subspace}{subruang}{a corridor inside a building: we may walk back and forth in that corridor, but we never leave it.}

\begin{untukanak}
Picture two straight corridors in a building that meet at a single point, like the letter X. Each
corridor, if we walk inside it, is still a corridor: we can go forward, backward, and stop
anywhere without leaving it. The same holds if we walk twice as far along the same corridor; we
are still in that corridor.

Now consider the union of the two corridors. Take one step in the first corridor, then one step in
the second. The place where we stop lies in neither corridor. So the union of two corridors is not
a corridor, even though each corridor is a corridor. This is the thing most often forgotten about
the word ``closed'': what is tested is not the members one by one, but the result of mixing them.

The nagging question: if three corridors meet at one point, is their union a corridor? Sometimes
yes, when all three are the same corridor. In this chapter we will learn why ``closed under
mixing'' is far stronger than ``contains the meeting point'', and why a single example of a step
that leaves is enough to disqualify everything.
\end{untukanak}

\section{Linear combinations and span}
\tingkat{intermediate}

A linear combination is the result of mixing several vectors in certain proportions. The span, in
turn, is the set of all those mixtures, for every conceivable proportion.

\begin{definisi}[Linear combination and span]
\index{linear combination}
\index{span}
A vector $\mathbf{v}$ is a \emph{linear combination} of
$\mathbf{v}_{1},\dots,\mathbf{v}_{k}$ if there are scalars $c_{1},\dots,c_{k}$ with
$\mathbf{v}=c_{1}\mathbf{v}_{1}+\dots+c_{k}\mathbf{v}_{k}$. The \emph{span} of a set
$S\subseteq V$, written $\operatorname{span}(S)$, is the set of all finite linear combinations
of members of $S$. For $S=\varnothing$ we stipulate
$\operatorname{span}(S)=\{\mathbf{0}\}$.
\end{definisi}

A span is always a subspace, and it is the smallest subspace containing $S$. The first statement
says the span deserves to be called a space; the second, that no other space is more economical.

\begin{lemma}
For every $S\subseteq V$, $\operatorname{span}(S)$ is a subspace of $V$, and every subspace
containing $S$ also contains $\operatorname{span}(S)$.
\end{lemma}

\par\medskip\noindent\textbf{Proof.}\ 
For $S=\varnothing$ the statement is true because $\{\mathbf{0}\}$ is a subspace. For $S$
nonempty, take $\mathbf{u}=\sum a_{i}\mathbf{s}_{i}$ and $\mathbf{v}=\sum b_{j}\mathbf{t}_{j}$,
two finite linear combinations of members of $S$. Then
$\mathbf{u}+\mathbf{v}=\sum a_{i}\mathbf{s}_{i}+\sum b_{j}\mathbf{t}_{j}$ and
$c\mathbf{u}=\sum (ca_{i})\mathbf{s}_{i}$ are also finite linear combinations of members of $S$,
so the span is closed under both operations; it also contains $\mathbf{0}=0\mathbf{s}$ for any
$\mathbf{s}\in S$. If $W$ is a subspace and $S\subseteq W$, then every $\mathbf{s}\in S$ lies in
$W$, and closure of $W$ forces every finite linear combination of members of $S$ to lie in $W$
as well.
\hfill$\square$\par\medskip

Therefore $\operatorname{span}(S)$ is the intersection of all subspaces containing $S$, and
adding a vector already inside the span changes nothing:
$\operatorname{span}(\operatorname{span}(S))=\operatorname{span}(S)$.

If $\mathbf{b}$ lies in the span of the columns of $A$, then the system $A\mathbf{x}=\mathbf{b}$
has a solution.

\begin{proposisi}
Let $A$ be an $m\times n$ matrix with columns $\mathbf{a}_{1},\dots,\mathbf{a}_{n}$ and let
$\mathbf{b}\in\mathbb{R}^{m}$. The system $A\mathbf{x}=\mathbf{b}$ is consistent if and only if
$\mathbf{b}\in\operatorname{span}(\mathbf{a}_{1},\dots,\mathbf{a}_{n})$.
\end{proposisi}

\par\medskip\noindent\textbf{Proof.}\ 
Writing $A\mathbf{x}=\mathbf{b}$ coordinate by coordinate is the same as writing
$x_{1}\mathbf{a}_{1}+\dots+x_{n}\mathbf{a}_{n}=\mathbf{b}$. So a solution $\mathbf{x}$ exists
exactly when $\mathbf{b}$ is a linear combination of the columns of $A$ with coefficients
$x_{1},\dots,x_{n}$.
\hfill$\square$\par\medskip

The set $\operatorname{span}(\mathbf{a}_{1},\dots,\mathbf{a}_{n})$ is called the \emph{column
space} of $A$, and it is a subspace of $\mathbb{R}^{m}$ by the lemma above. For the opening
table, that column space is built by four columns, and two of them are redundant.

As a small example, in $\mathbb{R}^{3}$ consider $\mathbf{s}_{1}=(1,0,1)$ and
$\mathbf{s}_{2}=(0,1,1)$. Every linear combination has the form $(a,b,a+b)$, so the span is the
plane $z=x+y$. The vector $(1,2,4)$ lies outside that plane because $4\neq 1+2$, whereas
$(1,2,3)$ lies inside it with $a=1$, $b=2$.

\begin{tangan}
Materials: two pencils of equal length and one flat table. Mark a point on the table with a
fingertip; call it the origin.

Place the first pencil with its tail at the origin. Place the second pencil with its tail at the
tip of the first, at a right angle to it. Now slide the second pencil without rotating it until
its tail is back at the origin. Its tip moves to another place, and that place is the sum of the
two pencils: two orders of addition stop at the same point.

Now the scalars. Reverse one pencil, and its direction becomes negative. Join two identical
pencils from the origin, and the length doubles. What the reader must see: all those tips stay on
one straight line, however far the pencils are joined. One pencil covers only one line on the
table; two nonparallel pencils cover the entire table surface. That is the difference between a
line and a plane, and no third pencil adds a new surface.
\end{tangan}



\section{Linear independence and dependence}
\tingkat{intermediate}

A set of vectors is called linearly independent if not one of its members can be mixed from the
other members. It is linearly dependent if some member is redundant.

\begin{definisi}[Linear independence and dependence]
\index{linear independence}
A set $\{\mathbf{v}_{1},\dots,\mathbf{v}_{k}\}\subseteq V$ is \emph{linearly independent} if the
only choice of scalars satisfying
$c_{1}\mathbf{v}_{1}+\dots+c_{k}\mathbf{v}_{k}=\mathbf{0}$ is
$c_{1}=\dots=c_{k}=0$. If some other choice exists, the set is \emph{linearly dependent}.
\end{definisi}

Two small conventions: the empty set is linearly independent because there are no coefficients to
check, and $\{\mathbf{0}\}$ is linearly dependent because $1\cdot\mathbf{0}=\mathbf{0}$ is a
nontrivial relation.

In a linearly dependent set there is always one vector that lies in the span of the preceding
vectors.

\begin{lemma}[Dependence]
Let $\mathbf{v}_{1}\neq\mathbf{0}$ and $S=\{\mathbf{v}_{1},\dots,\mathbf{v}_{k}\}$.
The set $S$ is linearly dependent if and only if there is an index $j\ge 2$ with
$\mathbf{v}_{j}\in\operatorname{span}(\mathbf{v}_{1},\dots,\mathbf{v}_{j-1})$.
\end{lemma}

\par\medskip\noindent\textbf{Proof.}\ 
If $\mathbf{v}_{j}=\sum_{i<j}c_{i}\mathbf{v}_{i}$, then
$\sum_{i<j}c_{i}\mathbf{v}_{i}-\mathbf{v}_{j}=\mathbf{0}$ is a nontrivial relation, because the
coefficient of $\mathbf{v}_{j}$ is $-1$; so $S$ is dependent. Conversely, if $S$ is dependent then
there is a nontrivial relation $\sum_{i=1}^{k}c_{i}\mathbf{v}_{i}=\mathbf{0}$. Choose the largest
$j$ with $c_{j}\neq 0$. Since $\mathbf{v}_{1}\neq\mathbf{0}$, the index $j$ is not $1$, so
$j\ge 2$ and
\[
  \mathbf{v}_{j}=-\frac{1}{c_{j}}\sum_{i<j}c_{i}\mathbf{v}_{i}
  \in\operatorname{span}(\mathbf{v}_{1},\dots,\mathbf{v}_{j-1}).
\]
\hfill$\square$\par\medskip

Three consequences we use constantly: a set containing $\mathbf{0}$ is always dependent; two
vectors are dependent exactly when one is a multiple of the other; and every subset of a linearly
independent set is also linearly independent.

As an example, consider $\mathbf{v}_{1}=(1,2,-1)$, $\mathbf{v}_{2}=(0,1,3)$, and
$\mathbf{v}_{3}=(2,5,1)$ in $\mathbb{R}^{3}$. Since
$2\mathbf{v}_{1}+\mathbf{v}_{2}=(2,5,1)=\mathbf{v}_{3}$, the relation
$2\mathbf{v}_{1}+\mathbf{v}_{2}-\mathbf{v}_{3}=\mathbf{0}$ is nontrivial and the three are
dependent. Not one of those vectors is zero, and no two are multiples; the dependence appears only
at three vectors.

Figure~\ref{fig:ch06-bidang} places the two side by side: the span of two vectors as a plane, and
three vectors one of which already lies in the span of the other two.

\begin{figure}[htbp]
\centering
\begin{tikzpicture}[scale=0.8,>=Stealth,font=\small]
  \begin{scope}[x={(1cm,0cm)}, y={(0.40cm,0.28cm)}, z={(0cm,1cm)}]
    \fill[layerteal!10] (0,0,0) -- (2.4,0,2.4) -- (2.4,1.8,4.2) -- (0,1.8,2.4) -- cycle;
    \draw[layerteal] (0,0,0) -- (2.4,0,2.4) -- (2.4,1.8,4.2) -- (0,1.8,2.4) -- cycle;
    \draw[->,layergrey] (0,0,0) -- (3.0,0,0) node[below] {$x$};
    \draw[->,layergrey] (0,0,0) -- (0,2.2,0) node[right] {$y$};
    \draw[->,layergrey] (0,0,0) -- (0,0,2.6) node[left] {$z$};
    \draw[->,thick,accent] (0,0,0) -- (1,0,1)
      node[above left,font=\scriptsize,inner sep=1pt] {$\mathbf{s}_{1}$};
    \draw[->,thick,layerindigo] (0,0,0) -- (0,1,1)
      node[right,font=\scriptsize,inner sep=1pt] {$\mathbf{s}_{2}$};
    \fill[layerjade] (1,2,3) circle (1.6pt);
    \node[right,font=\scriptsize] at (1,2,3) {$(1,2,3)$};
    \draw[ink] (1,2,4) circle (2pt);
    \node[right,font=\scriptsize] at (1,2,4) {$(1,2,4)$};
  \end{scope}
  \begin{scope}[xshift=6.4cm]
    \draw[->,layergrey] (-0.4,0) -- (3.2,0) node[below] {$x$};
    \draw[->,layergrey] (0,-0.4) -- (0,2.6) node[left] {$y$};
    \draw[dashed,layergrey] (0,0) -- (2,1);
    \draw[dashed,layergrey] (2,1) -- (2.4,2);
    \draw[->,thick,accent] (0,0) -- (1,0.5) node[below right] {$\mathbf{v}_{1}$};
    \draw[->,thick,layerindigo] (0,0) -- (0.4,1) node[above left] {$\mathbf{v}_{2}$};
    \draw[->,very thick,layerorange] (0,0) -- (2.4,2) node[right] {$\mathbf{v}_{3}$};
    \fill[ink] (0,0) circle (1.2pt);
  \end{scope}
\end{tikzpicture}
\caption{Left: $\operatorname{span}\{(1,0,1),(0,1,1)\}$ is the plane $z=x+y$; the point
$(1,2,3)$ lies inside it, the point $(1,2,4)$ outside it. Right: a two-dimensional sketch of
$\mathbf{v}_{1}=(1,2,-1)$ and $\mathbf{v}_{2}=(0,1,3)$ with
$\mathbf{v}_{3}=2\mathbf{v}_{1}+\mathbf{v}_{2}=(2,5,1)$; the dashed lines show
$2\mathbf{v}_{1}$ plus $\mathbf{v}_{2}$ landing exactly at the tip of $\mathbf{v}_{3}$, so that
the three are dependent.}
\label{fig:ch06-bidang}
\end{figure}



In $\mathbb{R}^{2}$ any three vectors are always dependent, without computing anything.

\begin{proposisi}
Every set of three vectors in $\mathbb{R}^{2}$ is linearly dependent.
\end{proposisi}

\par\medskip\noindent\textbf{Proof.}\ 
Let $S=\{\mathbf{u},\mathbf{v},\mathbf{w}\}\subseteq\mathbb{R}^{2}$. If two of them are multiples,
$S$ is dependent. If not, $\mathbf{u}$ and $\mathbf{v}$ are linearly independent. Take
$\mathbf{w}=(w_{1},w_{2})$ and look for $a,b$ with
$a\mathbf{u}+b\mathbf{v}=\mathbf{w}$. Writing out the components gives two linear equations in
$a$ and $b$, and since $\mathbf{u}$, $\mathbf{v}$ are not multiples, that system has a unique
solution. Hence $\mathbf{w}\in\operatorname{span}(\mathbf{u},\mathbf{v})$.
\hfill$\square$\par\medskip

In $\mathbb{R}^{3}$ a similar statement holds for four vectors, but the method above would become
a long case check; we need a uniform argument.

To test independence computationally, arrange the vectors as the columns of a matrix and then
solve $A\mathbf{c}=\mathbf{0}$; the set is linearly independent exactly when that system has only
the trivial solution. The determinant offers a shortcut in $\mathbb{R}^{3}$, but that tool is only
built in Chapter 8.

\section{Basis and dimension}
\tingkat{advanced}

A basis is a list of vectors that is enough to build the whole space and contains not one
redundant vector. The words ``enough'' and ``no redundancy'' are exactly the two conditions we
already have.

\begin{definisi}[Basis and dimension]
\index{basis}
A set $B\subseteq V$ is called a \emph{basis} of $V$ if $B$ is linearly independent and
$\operatorname{span}(B)=V$. If $V$ has a finite basis, $V$ is called
\emph{finite-dimensional}; if not, $V$ is infinite-dimensional.
\end{definisi}

The definition of dimension cannot yet be written, because it depends on a fact not yet proved:
two bases of the same space always have the same size.

First step: a representation in a basis is always unique.

\begin{lemma}[Uniqueness of coordinates]
If $B=\{\mathbf{b}_{1},\dots,\mathbf{b}_{n}\}$ is a basis of $V$, then every $\mathbf{v}\in V$
can be written uniquely as $\mathbf{v}=a_{1}\mathbf{b}_{1}+\dots+a_{n}\mathbf{b}_{n}$.
\end{lemma}

\par\medskip\noindent\textbf{Proof.}\ 
Existence of the representation comes from $\operatorname{span}(B)=V$. Suppose there are two
representations with coefficients $a_{i}$ and $a_{i}'$. Subtracting them gives
$\sum_{i}(a_{i}-a_{i}')\mathbf{b}_{i}=\mathbf{0}$. Because $B$ is linearly independent, all those
difference coefficients are zero, so $a_{i}=a_{i}'$ for every $i$.
\hfill$\square$\par\medskip

The numbers $a_{1},\dots,a_{n}$ are called the coordinates of $\mathbf{v}$ with respect to $B$;
change of basis is the subject of Chapter 7, and here we only need uniqueness.

Second step: one vector can be exchanged for another without destroying the spanning property.
Replacing one vector in a spanning set by a combination that contains it does not reduce what can
be built.

\begin{lemma}[Exchange]
\index{exchange lemma}
Let $\mathbf{u}=\sum_{\mathbf{s}\in S}a_{\mathbf{s}}\mathbf{s}$ be a finite sum with
$a_{\mathbf{s}_{0}}\neq 0$ for some $\mathbf{s}_{0}\in S$. If $S\cup\{\mathbf{u}\}$
spans $V$, then $(S\setminus\{\mathbf{s}_{0}\})\cup\{\mathbf{u}\}$ also spans $V$.
\end{lemma}

\par\medskip\noindent\textbf{Proof.}\ 
From that relation,
$\mathbf{s}_{0}=a_{\mathbf{s}_{0}}^{-1}\bigl(\mathbf{u}-\sum_{\mathbf{s}\neq\mathbf{s}_{0}}a_{\mathbf{s}}\mathbf{s}\bigr)$,
so $\mathbf{s}_{0}$ is a linear combination of members of the smaller set
$T=(S\setminus\{\mathbf{s}_{0}\})\cup\{\mathbf{u}\}$. Every member of $S\cup\{\mathbf{u}\}$
therefore lies in $\operatorname{span}(T)$, so
$\operatorname{span}(S\cup\{\mathbf{u}\})\subseteq\operatorname{span}(T)$. The reverse inclusion
is clear because $T\subseteq S\cup\{\mathbf{u}\}$. Hence both spans equal $V$.
\hfill$\square$\par\medskip

The third step is the heart of this chapter.

If a space can be built by $n$ vectors, then it is impossible to find $n+1$ linearly independent
vectors inside it.

\begin{teorema}[Exchange]
\index{exchange theorem}
If $V$ is spanned by $n$ vectors and $\mathbf{w}_{1},\dots,\mathbf{w}_{m}$ are linearly
independent in $V$, then $m\le n$.
\end{teorema}

\par\medskip\noindent\textbf{Proof.}\ 
Let $S_{0}=\{\mathbf{v}_{1},\dots,\mathbf{v}_{n}\}$ span $V$. For $j=1$:
$\mathbf{w}_{1}\neq\mathbf{0}$ is a linear combination of members of $S_{0}$ with at least one
nonzero coefficient, so the exchange lemma gives $T_{1}\subseteq S_{0}$ of size $n-1$ with
$\{\mathbf{w}_{1}\}\cup T_{1}$ spanning.

Suppose $\{\mathbf{w}_{1},\dots,\mathbf{w}_{j}\}\cup T_{j}$ spans with $|T_{j}|=n-j$. The vector
$\mathbf{w}_{j+1}$ is a linear combination of that set. If all coefficients on $T_{j}$ are zero,
then $\mathbf{w}_{j+1}$ becomes a linear combination of $\mathbf{w}_{1},\dots,\mathbf{w}_{j}$,
which is impossible because that set is linearly independent. So some member of $T_{j}$ has a
nonzero coefficient, and the exchange lemma gives $T_{j+1}$ of size $n-j-1$ that still spans
together with $\mathbf{w}_{1},\dots,\mathbf{w}_{j+1}$.

After $m$ steps we obtain a spanning set of size $n-m$; the size of a set cannot be negative, so
$m\le n$.
\hfill$\square$\par\medskip

Now the size of a space can be written: the number of vectors in its most economical list, and
that number does not depend on the list chosen.

Two bases of a finite-dimensional space always have the same number of elements.

\begin{akibat}
If $V$ is finite-dimensional with a basis $B_{1}$ of size $n$ and a basis $B_{2}$ of size $m$,
then $n=m$.
\end{akibat}

\par\medskip\noindent\textbf{Proof.}\ 
$B_{1}$ spans and $B_{2}$ is linearly independent, so the exchange theorem gives $m\le n$.
Conversely $B_{2}$ spans and $B_{1}$ is linearly independent, hence $n\le m$.
\hfill$\square$\par\medskip

The number of vectors in that most economical list is important enough to deserve a name of its
own.

\begin{definisi}[Dimension]
\index{dimension}
For finite-dimensional $V$, $\dim V$ is the number of members of any basis of $V$. The null space
$\{\mathbf{0}\}$ is stipulated to have dimension $0$.
\end{definisi}

The size of a subspace never exceeds the size of its space; a list whose length matches the
dimension is immediately a basis; and a linearly independent list can always be completed to full
length.

\begin{akibat}
Let $\dim V=n$. (a) Every linearly independent set of $n$ vectors is a basis.
(b) Every spanning set of $n$ vectors is a basis. (c) Every linearly independent set
can be extended to a basis. (d) If $U$ is a subspace of $V$, then $\dim U\le\dim V$, and
equality holds only when $U=V$.
\end{akibat}

\par\medskip\noindent\textbf{Proof.}\ 
For (a), a linearly independent set of size $n$ that does not span can be enlarged by one vector
outside its span, and the result is still linearly independent; so there are $n+1$ linearly
independent vectors, and the exchange theorem forbids that. For (b), a spanning set of size $n$
that is dependent can lose one vector through the dependence lemma and still span with $n-1$
vectors; again this contradicts exchange. For (c), add a vector outside the span as long as one
exists; the process stops after at most $n$ steps, and when it stops the span is already all of
$V$. For (d), a basis of $U$ is a linearly independent set in $V$, so its size is at most $n$; if
its size is $n$, part (a) makes it a basis of $V$, hence $U=V$.
\hfill$\square$\par\medskip

Thus $\dim\mathbb{R}^{n}=n$, $\dim P_{n}=n+1$ with basis $1,x,\dots,x^{n}$,
$\dim M_{m\times n}(\mathbb{R})=mn$ with basis the matrices having exactly one entry equal to $1$,
and $\dim\{\mathbf{0}\}=0$. For function spaces such as $\mathbb{R}^{\mathbb{R}}$ there is no
finite basis; the proof requires the axiom of choice and lies outside the scope of this book.

\kamus{basis}{basis}{the most economical list of paint cans: enough to make every color, and not one can in it can be thrown away.}

\kamus{dimension}{dimensi}{the number of cans in that most economical list; the number belongs to the space, not to the list we happen to choose.}

\section{Direct sums and complementary subspaces}
\tingkat{advanced}

The sum of two subspaces is the set of all results of adding one member of each subspace.

\begin{definisi}[Sum of subspaces]
\index{sum of subspaces}
For subspaces $U,W\subseteq V$, the \emph{sum} $U+W$ is
$\{\mathbf{u}+\mathbf{w}:\mathbf{u}\in U,\ \mathbf{w}\in W\}$.
\end{definisi}

It differs from the union, which is almost never a subspace, as the two-axes example shows.

The sum is itself a subspace, and no smaller subspace contains both.

\begin{proposisi}
For subspaces $U,W\subseteq V$, the set $U+W$ is a subspace of $V$, and it is contained in every
subspace that contains $U\cup W$.
\end{proposisi}

\par\medskip\noindent\textbf{Proof.}\ 
The set $U+W$ contains $\mathbf{0}=\mathbf{0}+\mathbf{0}$. If
$\mathbf{u}_{1}+\mathbf{w}_{1}$ and $\mathbf{u}_{2}+\mathbf{w}_{2}$ are two of its members, their
sum $(\mathbf{u}_{1}+\mathbf{u}_{2})+(\mathbf{w}_{1}+\mathbf{w}_{2})$ has its first term in $U$
and its second in $W$, and $c(\mathbf{u}_{1}+\mathbf{w}_{1})=c\mathbf{u}_{1}+c\mathbf{w}_{1}$ is
of the same form; so $U+W$ is a subspace. If $Z$ is a subspace containing $U$ and $W$, then $Z$
contains every $\mathbf{u}+\mathbf{w}$ because it is closed under addition, so $U+W\subseteq Z$.
\hfill$\square$\par\medskip

The size of a sum is not the sum of the sizes, because the part that overlaps is counted twice.

$\dim(U+W)$ equals $\dim U$ plus $\dim W$ minus the dimension of the intersection.

\begin{teorema}[Dimension of a sum]
\index{dimension of a sum}
If $U$ and $W$ are subspaces of a finite-dimensional space $V$, then
\[
  \dim(U+W)=\dim U+\dim W-\dim(U\cap W).
\]
\end{teorema}

\par\medskip\noindent\textbf{Proof.}\ 
The intersection $U\cap W$ is a subspace, so it is finite-dimensional. Take a basis
$\{\mathbf{z}_{1},\dots,\mathbf{z}_{m}\}$ of $U\cap W$. Since it is linearly independent in $U$,
that set is extended to a basis of $U$ by adding $\mathbf{u}_{1},\dots,
\mathbf{u}_{j}$, so that $\dim U=m+j$; in the same way, adding
$\mathbf{w}_{1},\dots,\mathbf{w}_{k}$ yields a basis of $W$, so that $\dim W=m+k$.

Claim: $\{\mathbf{z}_{1},\dots,\mathbf{z}_{m},\mathbf{u}_{1},\dots,\mathbf{u}_{j},
\mathbf{w}_{1},\dots,\mathbf{w}_{k}\}$ is a basis of $U+W$. The set contains a basis of $U$ and a
basis of $W$, so its span contains $U+W$; the reverse inclusion is clear. For independence,
suppose
\[
  \sum_{i}a_{i}\mathbf{z}_{i}+\sum_{r}b_{r}\mathbf{u}_{r}
  +\sum_{s}c_{s}\mathbf{w}_{s}=\mathbf{0}.
\]
Then $\sum_{s}c_{s}\mathbf{w}_{s}=-\sum_{i}a_{i}\mathbf{z}_{i}-\sum_{r}b_{r}\mathbf{u}_{r}$
lies in $U$ and at the same time in $W$; so it is a member of $U\cap W$ and can be written
$\sum_{i}d_{i}\mathbf{z}_{i}$. Then $\sum_{s}c_{s}\mathbf{w}_{s}-\sum_{i}d_{i}\mathbf{z}_{i}
=\mathbf{0}$ is a linear relation in the basis of $W$, so all $c_{s}$ and $d_{i}$ are zero; the
original equation shrinks to a linear relation in the basis of $U$, so that $a_{i}$ and $b_{r}$
are also zero.

Hence $\dim(U+W)=m+j+k=(m+j)+(m+k)-m=\dim U+\dim W-\dim(U\cap W)$.
\hfill$\square$\par\medskip

Now the second question from the opening of the chapter: when is the decomposition into two terms
unique?

The sum $U+W$ is called a \emph{direct sum}, written $U\oplus W$, if every member of
$U+W$ can be written as $\mathbf{u}+\mathbf{w}$ in exactly one way.

\begin{definisi}[Direct sum]
\index{direct sum}
If every $\mathbf{x}\in U+W$ has exactly one representation
$\mathbf{x}=\mathbf{u}+\mathbf{w}$ with $\mathbf{u}\in U$ and $\mathbf{w}\in W$, then
$U+W$ is called a direct sum and is written $U\oplus W$.
\end{definisi}

Uniqueness of all representations can be checked by looking only at the representation of zero.

$U+W$ is a direct sum exactly when $U\cap W=\{\mathbf{0}\}$.

\begin{teorema}[Criterion for a direct sum]
Let $U$ and $W$ be subspaces of $V$. Then $U+W=U\oplus W$ if and only if
$U\cap W=\{\mathbf{0}\}$.
\end{teorema}

\par\medskip\noindent\textbf{Proof.}\ 
Suppose $U+W$ is a direct sum and take $\mathbf{x}\in U\cap W$. Then
$\mathbf{0}=\mathbf{x}+(-\mathbf{x})$ with $\mathbf{x}\in U$ and $-\mathbf{x}\in W$, and also
$\mathbf{0}=\mathbf{0}+\mathbf{0}$. Uniqueness gives $\mathbf{x}=\mathbf{0}$.

Conversely, suppose $U\cap W=\{\mathbf{0}\}$. Take $\mathbf{x}\in U+W$ with two representations
$\mathbf{x}=\mathbf{u}_{1}+\mathbf{w}_{1}=\mathbf{u}_{2}+\mathbf{w}_{2}$. Then
$\mathbf{u}_{1}-\mathbf{u}_{2}=\mathbf{w}_{2}-\mathbf{w}_{1}$ lies in $U$ and in $W$, so it equals
$\mathbf{0}$. Consequently $\mathbf{u}_{1}=\mathbf{u}_{2}$ and $\mathbf{w}_{1}=\mathbf{w}_{2}$.
\hfill$\square$\par\medskip

For two subspaces that criterion is enough. For three subspaces or more, checking pairwise
intersections is not enough, and the next example shows why.

In $\mathbb{R}^{3}$ take $V_{1}=\operatorname{span}\{(1,0,0)\}$,
$V_{2}=\operatorname{span}\{(0,1,0)\}$, and $V_{3}=\operatorname{span}\{(1,1,0)\}$. Every
pairwise intersection contains only $\mathbf{0}$, but $V_{1}+V_{2}+V_{3}$ is not a direct sum:
$(1,1,0)$ can be written $(1,0,0)+(0,1,0)+\mathbf{0}$ and also $\mathbf{0}+\mathbf{0}+(1,1,0)$.
Even zero has two representations, namely $\mathbf{0}+\mathbf{0}+\mathbf{0}$ and
$(1,0,0)+(0,1,0)+(-1,-1,0)$.

For many terms, the criterion returns to the definition:
$\mathbf{0}=\mathbf{v}_{1}+\dots+\mathbf{v}_{k}$ with $\mathbf{v}_{i}\in V_{i}$ is possible only
when all the $\mathbf{v}_{i}$ are zero.

Finally, every subspace can be paired with another subspace so that the two decompose the whole
space.

\begin{teorema}[Complement]
\index{complement}
If $V$ is finite-dimensional and $U$ is a subspace of $V$, then there is a subspace
$W\subseteq V$ with $V=U\oplus W$. In that case $\dim V=\dim U+\dim W$.
\end{teorema}

\par\medskip\noindent\textbf{Proof.}\ 
If $U=\{\mathbf{0}\}$, take $W=V$. If not, take a basis
$\{\mathbf{u}_{1},\dots,\mathbf{u}_{m}\}$ of $U$ and extend it to a basis
$\{\mathbf{u}_{1},\dots,\mathbf{u}_{m},\mathbf{w}_{1},\dots,\mathbf{w}_{k}\}$ of $V$;
this extension is possible by consequence (c) in the previous section. Set
$W=\operatorname{span}\{\mathbf{w}_{1},\dots,\mathbf{w}_{k}\}$.

Every $\mathbf{x}\in V$ is a linear combination of the entire basis, and separating the
$\mathbf{u}$ terms from the $\mathbf{w}$ terms exhibits it as a member of $U+W$. For the
intersection, take $\mathbf{x}\in U\cap W$. Then
$\mathbf{x}=\sum a_{i}\mathbf{u}_{i}=\sum b_{s}\mathbf{w}_{s}$, so that
$\sum a_{i}\mathbf{u}_{i}-\sum b_{s}\mathbf{w}_{s}=\mathbf{0}$. This is a linear relation in the
basis of $V$, so all its coefficients are zero, in particular $\mathbf{x}=\mathbf{0}$. By the
criterion above, $V=U\oplus W$, and from the dimension-of-a-sum theorem,
$\dim V=\dim U+\dim W$.

The complement is not unique. In $\mathbb{R}^{2}$ with $U$ the $x$-axis, every line through the
origin other than the $x$-axis is a complement of $U$; what is unique is not the choice but the
dimension. In Chapter 10, orthogonal complements make that decomposition the shortest one as well.

Now the opening question can be answered. For $\mathbf{w}=(1,1,1)$, let
$U=\{\boldsymbol{\delta}:\mathbf{w}\cdot\boldsymbol{\delta}=0\}$ and
$W=\operatorname{span}\{\mathbf{w}\}$. Their intersection is only $\{\mathbf{0}\}$, since
$\mathbf{w}\cdot(t\mathbf{w})=3t$ is zero only when $t=0$; with $\dim U=2$ and
$\dim W=1$ we obtain $U\oplus W=\mathbb{R}^{3}$. Every perturbation splits uniquely into a part
that does not change the score, namely a member of $U$, and a part that changes it, a member of
$W$.

\kamus{direct sum}{jumlah langsung}{two corridors that meet only at the origin, so that every place in the building is reachable by one route that cannot be disputed.}

\section{Worked examples}
\tingkat{intermediate}

The following four computations use the entire content of the chapter; all of them are rechecked
with Python in the lab section.

\subsection*{(a) Rank and a basis of the column space}

Back to the opening table. Row elimination: $R_{2}\leftarrow R_{2}-2R_{1}$ turns the second row
into $(0,0,1,1)$, and $R_{3}\leftarrow R_{3}-3R_{1}$ turns the third row into
$(0,0,2,2)$. Then $R_{3}\leftarrow R_{3}-2R_{2}$ eliminates the third row, so that what remains
is $\mathbf{r}_{1}=(1,2,0,1)$ and $(0,0,1,1)$.

That form is already reduced echelon, with pivot columns $1$ and $3$. Hence
$\operatorname{rank}A=2$ and its column space has dimension $2$ with basis
$\{\mathbf{c}_{1},\mathbf{c}_{3}\}=\{(1,2,3),(0,1,2)\}$.

A direct check against the discarded columns:
\[
  \mathbf{c}_{2}=(2,4,6)=2\mathbf{c}_{1},
  \qquad
  \mathbf{c}_{4}=(1,3,5)=\mathbf{c}_{1}+\mathbf{c}_{3}.
\]
Both are exact, with no rounding. This is the answer ``two'' to the opening question: the pivot
columns give a linearly independent set that spans the column space, and a column without a pivot
is a combination of the pivot columns to its left.

The null space of $A$ is also readable from the reduced form. From $x_{1}+2x_{2}+x_{4}=0$ and
$x_{3}+x_{4}=0$ we obtain $x_{1}=-2x_{2}-x_{4}$ and $x_{3}=-x_{4}$, so that
\[
  \mathbf{x}=x_{2}\,(-2,1,0,0)+x_{4}\,(-1,0,-1,1),
\]
written as row vectors for brevity. Hence $\dim\ker A=2$, and
$\operatorname{rank}A+\dim\ker A=2+2=4$ equals the number of columns; the general statement is
rank-nullity, proved in Chapter 7.

\subsection*{(b) Three dependent vectors}

Consider $\mathbf{v}_{1}=(1,2,-1)$, $\mathbf{v}_{2}=(0,1,3)$, $\mathbf{v}_{3}=(2,5,1)$ as
the columns of $V$. Cofactor expansion along the first row:
\[
  \det V=1(1\cdot1-3\cdot5)+0+2\bigl(2\cdot3-1\cdot(-1)\bigr)=(1-15)+2(6+1)=0.
\]
A zero determinant means the columns are dependent; the relation is guessed and then checked:
$2\mathbf{v}_{1}+\mathbf{v}_{2}=(2,5,1)=\mathbf{v}_{3}$, so
$2\mathbf{v}_{1}+\mathbf{v}_{2}-\mathbf{v}_{3}=\mathbf{0}$ and $\operatorname{rank}V=2$.

\subsection*{(c) Dimension of a polynomial space}

$\dim P_{3}=4$, because $\{1,x,x^{2},x^{3}\}$ is a basis. The set
$B=\{1,\ 1+x,\ 1+x+x^{2},\ 1+x+x^{2}+x^{3}\}$ is also a basis: its coordinate matrix with respect
to the standard basis is upper triangular with diagonal $1$, so $B$ is linearly independent, and
since its size is $4=\dim P_{3}$, consequence (a) makes it a basis.

The coordinates of $p=2-x+3x^{3}$ with respect to $B$ are computed from the highest degree
downward: $a_{4}=3$, then $a_{3}+a_{4}=0$ gives $a_{3}=-3$, $a_{2}+a_{3}+a_{4}=-1$ gives
$a_{2}=-1$, and the constant term gives $a_{1}=3$. Hence
$p=3-(1+x)-3(1+x+x^{2})+3(1+x+x^{2}+x^{3})$, which upon expansion gives $2-x+3x^{3}$.

\subsection*{(d) Decomposition as a direct sum}

Take $U=\{(x,y,z):x+y+z=0\}$ and $W=\operatorname{span}\{(1,1,1)\}$ in $\mathbb{R}^{3}$.
A basis of $U$ is $\{(1,-1,0),(1,0,-1)\}$, so $\dim U=2$; with $\dim W=1$ and
$U\cap W=\{\mathbf{0}\}$ we obtain $\mathbb{R}^{3}=U\oplus W$.

Decompose $\mathbf{v}=(3,0,1)$ as $\mathbf{u}+t(1,1,1)$. The condition that the sum of the
coordinates of $\mathbf{u}$ be zero gives $3+0+1-3t=0$, so $t=4/3$. Then
\[
  \mathbf{v}=\underbrace{\Bigl(\tfrac53,-\tfrac43,-\tfrac13\Bigr)}_{\in U}
  +\underbrace{\Bigl(\tfrac43,\tfrac43,\tfrac43\Bigr)}_{\in W},
\]
and $\tfrac53-\tfrac43-\tfrac13=0$ confirms that the first term really lies in $U$. By the
criterion theorem, there is no other decomposition.

\kertas{Redo (a) without looking: carry out the elimination on $A$, mark the pivot columns, then
write $\mathbf{c}_{2}$ and $\mathbf{c}_{4}$ as combinations of $\mathbf{c}_{1}$ and
$\mathbf{c}_{3}$. Then compute $\dim\ker A$ from the reduced form and check that its sum with the
rank equals the number of columns. Finally, for the matrix $V$ in (b), replace entry $(3,1)$ by
$0$ and decide whether the columns are still dependent.}

\section{Computational lab}
\tingkat{advanced}

The following code runs as is with numpy. The function \texttt{rref} is written by hand so that the
pivot columns are visible; its last line cleans up leftover rounding and negative zeros.

\begin{lab}
\begin{verbatim}
import numpy as np

np.set_printoptions(precision=3, suppress=True)

def rref(M):
    M = M.astype(float).copy()
    r = 0
    for c in range(M.shape[1]):
        i = np.argmax(np.abs(M[r:, c])) + r
        if abs(M[i, c]) < 1e-12:
            continue
        M[[r, i]] = M[[i, r]]
        M[r] = M[r] / M[r, c]
        for j in range(M.shape[0]):
            if j != r:
                M[j] = M[j] - M[j, c] * M[r]
        r = r + 1
        if r == M.shape[0]:
            break
    M[np.abs(M) < 1e-12] = 0.0
    return M

A = np.array([[1, 2, 0, 1],
              [2, 4, 1, 3],
              [3, 6, 2, 5]], float)
print("RREF(A):"); print(rref(A))
print("rank(A) =", np.linalg.matrix_rank(A))
print("k2-2k1 =", A[:, 1] - 2 * A[:, 0])
print("k4-k1-k3 =", A[:, 3] - A[:, 0] - A[:, 2])

V = np.array([[1, 0, 2],
              [2, 1, 5],
              [-1, 3, 1]], float)
print("det(V) =", np.linalg.det(V))
print("2v1+v2-v3 =", 2 * V[:, 0] + V[:, 1] - V[:, 2])

x = np.array([0., 1, 2, 3])
M = np.vander(x, 4, increasing=True)
print("det(M) =", np.linalg.det(M))
print("rank(M) =", np.linalg.matrix_rank(M))

B = np.array([[1, 1, 1, 1],
              [0, 1, 1, 1],
              [0, 0, 1, 1],
              [0, 0, 0, 1]], float)
p = np.array([2, -1, 0, 3], float)
k = np.linalg.solve(B, p)
print("koordinat p =", k)
print("sisa =", B @ k - p)
\end{verbatim}
\end{lab}

The printed output:
\begin{verbatim}
RREF(A):
[[1. 2. 0. 1.]
 [0. 0. 1. 1.]
 [0. 0. 0. 0.]]
rank(A) = 2
k2-2k1 = [0. 0. 0.]
k4-k1-k3 = [0. 0. 0.]
det(V) = 0.0
2v1+v2-v3 = [0. 0. 0.]
det(M) = 12.0
rank(M) = 4
koordinat p = [ 3. -1. -3.  3.]
sisa = [0. 0. 0. 0.]
\end{verbatim}

Four things are readable from that output. The pivot columns $1$ and $3$ as well as
$\operatorname{rank}A=2$ agree with the hand computation. The differences
$\mathbf{c}_{2}-2\mathbf{c}_{1}$ and $\mathbf{c}_{4}-\mathbf{c}_{1}-\mathbf{c}_{3}$ are exactly
zero; on real data those differences are small but nonzero, and that is where the limits of
applying rank numerically lie. $\det V=0$ confirms the dependence in (b), and the Vandermonde
matrix of the points $0,1,2,3$ has determinant $12$ and rank $4$. The coordinates of $p$ are
printed as $(3,-1,-3,3)$, matching the hand computation.

\begin{untukdosen}
Start from the opening table, not from the axioms: ask the students to guess how many columns are
independent, then show that the guess ``four'' is wrong. The axioms come in after that, as an
answer to a need. The subspace test and span in the second meeting, linear independence and basis
in the third, direct sums in the fourth.

Two typical misconceptions. First, testing a subspace only by ``does it contain zero'' and then
stopping; the first quadrant contains zero and is closed under addition, but fails on negative
multiples. Second, taking basis to mean standard basis, so that $(5,7)$ and $(4,3)$ are judged not
to be a basis of $\mathbb{R}^{2}$; stress that being a basis is a property of the set.

Class activity: each group writes three vectors in $\mathbb{R}^{3}$ and decides independence or
dependence by elimination; another group changes one entry and asks for a new decision. Students
are ready to move on when they can explain why three vectors in $\mathbb{R}^{2}$ are always
dependent.
\end{untukdosen}

\section{Summary}
\tingkat{basic}

\begin{ringkasan}
\begin{itemize}
  \item Vector space: a set with addition and scalar multiplication satisfying the eight
        axioms; its members may be vectors, matrices, polynomials, or functions.
  \item A subspace is checked with just three things: it contains $\mathbf{0}$ and is closed
        under both operations.
  \item $\operatorname{span}(S)$ is the smallest subspace containing $S$; $\mathbf{b}$ lies in
        the column space of $A$ exactly when $A\mathbf{x}=\mathbf{b}$ is consistent.
  \item Linear dependence: one member lies in the span of the preceding members; three vectors
        in $\mathbb{R}^{2}$ are always dependent.
  \item The exchange theorem traps $n+1$ linearly independent vectors in a space spanned by $n$
        vectors, so that all bases have the same size and $\dim V$ is defined.
  \item $\dim(U+W)=\dim U+\dim W-\dim(U\cap W)$; $U+W$ is direct exactly when
        $U\cap W=\{\mathbf{0}\}$; every finite-dimensional subspace has a complement.
\end{itemize}
\end{ringkasan}

\section{Exercises}
\tingkat{basic}

\begin{latihan}
\begin{enumerate}
  \item Decide whether $W=\{(x,y,z)\in\mathbb{R}^{3}:x-2y+z=0\ \text{and}\ x+y=0\}$
        is a subspace of $\mathbb{R}^{3}$. If so, find a basis and its dimension.
  \item Prove that the set of all $2\times2$ matrices whose main diagonal entries are zero
        is a subspace of $M_{2\times2}(\mathbb{R})$, then find a basis and its dimension.
  \item Show that $S=\{1+x,\ 1+x^{2},\ 2+x+x^{2}\}$ is linearly dependent in $P_{2}$.
        Find $\dim\operatorname{span}(S)$, then extend a largest linearly independent
        subset of $S$ to a basis of $P_{2}$.
  \item Prove: if $\{\mathbf{v}_{1},\mathbf{v}_{2},\mathbf{v}_{3}\}$ is linearly independent in
        $V$, then $\{\mathbf{v}_{1}+\mathbf{v}_{2},\ \mathbf{v}_{2}+\mathbf{v}_{3},\
        \mathbf{v}_{3}+\mathbf{v}_{1}\}$ is also linearly independent.
  \item Compute $\dim(U\cap W)$ and $\dim(U+W)$ for the two subspaces
        \begin{align*}
          U &= \operatorname{span}\{(1,0,0),(0,1,0)\}, \\
          W &= \operatorname{span}\{(1,1,0),(0,0,1)\}
        \end{align*}
        in $\mathbb{R}^{3}$. Check the dimension-of-a-sum formula, then compute them with numpy
        as well.
  \item For the $3\times4$ matrix with rows $(1,2,3,4)$, $(2,4,6,9)$, and $(0,1,1,1)$,
        find the rank, one basis of the column space, and one basis of the null space.
        Check with the function \texttt{matrix\_rank} from numpy and
        \texttt{null\_space} from scipy.
  \item Assess this claim: \emph{if $U_{1},U_{2},U_{3}$ are subspaces of $V$ and every pairwise
        intersection equals $\{\mathbf{0}\}$, then
        $U_{1}+U_{2}+U_{3}=U_{1}\oplus U_{2}\oplus U_{3}$.} State true, false, or not enough
        evidence, and include a reason or a counterexample.
  \item Prove that if $\dim V=n$ and $U$ is a subspace with $\dim U=n-1$, then there is a
        vector $\mathbf{v}\in V$ such that $V=U\oplus\operatorname{span}\{\mathbf{v}\}$.
  \item Find the dimension of the space of all polynomials of degree at most $4$ that vanish at
        $x=0$ and $x=1$ simultaneously, then give one basis of it. Check your answer by
        computing the rank of the $2\times5$ evaluation matrix at those two points.
  \item Given $B=\{1,\ 1+x,\ 1+x+x^{2},\ 1+x+x^{2}+x^{3}\}$ a basis of $P_{3}$. Find the
        coordinates of $q=4+2x-5x^{2}+x^{3}$ with respect to $B$, then check by solving the
        linear system in numpy.
\end{enumerate}

\noindent\textbf{Hints for the three hardest items.} Item 4: suppose
$\sum c_{i}(\dots)=\mathbf{0}$, collect the coefficient of each $\mathbf{v}_{i}$, then use the
independence of $\{\mathbf{v}_{1},\mathbf{v}_{2},\mathbf{v}_{3}\}$. Item 8: take a basis of $U$,
extend it by one vector, then use consequence (d). Item 9: a polynomial of degree at most
$4$ that vanishes at $0$ and $1$ is divisible by $x(x-1)$; count the free parameters.
\end{latihan}

\begin{untukpeneliti}
Three things are genuinely proved in this chapter, and all three hold over any field: the
uniqueness of the size of a basis through the exchange theorem (Steinitz's result), the formula
$\dim(U+W)=\dim U+\dim W-\dim(U\cap W)$, and the existence of a complement for every
finite-dimensional subspace. The popular claim that ``rank equals the number of informative
features'' is meaningful only if ``informative'' means linearly independent, and on real data that
condition is never satisfied exactly.

Primary sources: Anton, \emph{Elementary Linear Algebra} 11th edition, Chapter 4 sections
4.1--4.5; Axler, \emph{Linear Algebra Done Right} 4th edition, sections 1B, 1C, 2A--2C; Hoffman and
Kunze, \emph{Linear Algebra} 2nd edition, Chapter 2.

Open question: how far does $\dim(U\cap W)$ shift as a result of perturbing the entries? Limits of
validity: all results assume the vector space axioms; for infinite-dimensional spaces, the
existence of a basis requires Zorn's lemma. A rank computed by a computer depends on a threshold.
\end{untukpeneliti}

\peta{This chapter follows Anton, \emph{Elementary Linear Algebra} 11th edition, Chapter 4
(General Vector Spaces) section 4.1 (Real Vector Spaces) for the axioms and examples, 4.2
(Subspaces) for the subspace test and span, 4.3 (Linear Independence) for independence and the
dependence lemma, and 4.5 (Dimension) for the exchange theorem and its consequences; other
editions number these differently. A more abstract counterpart is in Axler, \emph{Linear Algebra
Done Right} 4th edition, Chapter 1 sections 1B and 1C (including sums and direct sums of
subspaces) and Chapter 2 sections 2A, 2B, and 2C, where the exchange theorem appears as a
comparison of the length of a linearly independent list with that of a spanning list. The
treatment of coordinates, linear transformations, and rank-nullity is not in this chapter but in
Chapter 7.}

